\documentclass[11pt,a4paper]{article}
\usepackage[utf8]{inputenc}
\usepackage{amsmath,amssymb,amsthm}
\usepackage{booktabs}
\usepackage{longtable}
\usepackage[dvipsnames]{xcolor}
\usepackage{tcolorbox}
\usepackage[margin=2.5cm]{geometry}
\usepackage{hyperref}
\usepackage{enumitem}

\tcbuselibrary{skins,breakable}

\newtcolorbox{beforebox}{colback=red!5,colframe=red!75!black,title=BEFORE (Original Claim),fonttitle=\bfseries,breakable}
\newtcolorbox{afterbox}{colback=green!5,colframe=green!50!black,title=AFTER (Salvageable Version),fonttitle=\bfseries,breakable}
\newtcolorbox{whywrongbox}{colback=orange!5,colframe=orange!75!black,title=WHY WRONG,fonttitle=\bfseries,breakable}
\newtcolorbox{insightbox}{colback=blue!5,colframe=blue!50!black,title=KEY INSIGHT,fonttitle=\bfseries,breakable}
\newtcolorbox{warningbox}{colback=yellow!5,colframe=yellow!75!black,title=WARNING,fonttitle=\bfseries,breakable}

\newtheorem{theorem}{Theorem}
\newtheorem{proposition}{Proposition}
\newtheorem{corollary}{Corollary}
\newtheorem{remark}{Remark}

\title{Proof Audit Report:\\Streaming Mirror-LoRA Theoretical Claims}
\author{Cross-Model Adversarial Review\\GPT-5.4 (xhigh) + Claude Opus 4}
\date{April 14, 2026 --- Round 1}

\begin{document}
\maketitle

\begin{abstract}
This report documents the rigorous verification of five theoretical claims
(3~theorems + 3~corollaries) in the paper
\emph{Dual-Timescale Adaptation Is Necessary and Sufficient for Streaming
Low-Rank Learning under Multi-Rate Drift}.
Cross-model adversarial review identified \textbf{30 issues} including
\textbf{8~FATAL}, \textbf{10~CRITICAL}, \textbf{7~MAJOR}, and
\textbf{5~MINOR} problems.
The acceptance gate is \textbf{NOT MET}.
This report details each issue with before/after analysis and salvage recommendations.
\end{abstract}

\tableofcontents
\newpage

%% ======================================================================
\section{Overview}
%% ======================================================================

\begin{center}
\begin{tabular}{lcccc}
\toprule
\textbf{Severity} & \textbf{Count} & \textbf{Fixed} & \textbf{Open} & \textbf{Gate Req.} \\
\midrule
FATAL    & 8  & 0 & 8  & 0 required \\
CRITICAL & 10 & 0 & 10 & 0 required \\
MAJOR    & 7  & 0 & 7  & --- \\
MINOR    & 5  & 0 & 5  & --- \\
\midrule
\textbf{Total} & \textbf{30} & \textbf{0} & \textbf{30} & \\
\bottomrule
\end{tabular}
\end{center}

\begin{warningbox}
\textbf{No formal proofs exist anywhere in the repository.}
All theoretical claims are one-line assertions in proposal documents.
The synthetic benchmark provides empirical illustrations, not proofs.
\end{warningbox}

%% ======================================================================
\section{FATAL Issue \#2: Strong Convexity Assumption}
%% ======================================================================

\begin{beforebox}
\textbf{Claim (FINAL\_PROPOSAL.md:33):}
``Streaming prediction with loss $\ell_t(\theta)$ that is $\mu$-strongly convex
and $L$-smooth in adapted low-rank parameters.''

Used by: Theorem~1 (Necessity), Theorem~2 (Sufficiency), all corollaries.
\end{beforebox}

\begin{whywrongbox}
\textbf{Counterexample (verified):}
The per-sample squared loss
$\ell_t(\theta) = \tfrac{1}{2}(\theta^\top x_t - y_t)^2$
has Hessian $H_t = x_t x_t^\top$, which is rank-1 for $d > 1$.
This is \emph{not} strongly convex. Even the project's own synthetic
benchmark (\texttt{synthetic\_benchmark.py:129}) uses this loss.

For LLM NTP loss, the landscape is highly non-convex even restricted
to the LoRA subspace. Strong convexity is indefensible without heavy
additional assumptions (e.g., batch loss with well-conditioned features).
\end{whywrongbox}

\begin{afterbox}
\textbf{Salvageable version:}
\begin{itemize}[nosep]
\item Use the \emph{expected} (batch) loss $L(\theta) = \mathbb{E}[\ell_t(\theta)]$,
      which for Gaussian $x_t \sim \mathcal{N}(0, \Sigma)$ gives
      $\nabla^2 L = \Sigma$, PD if $\Sigma \succ 0$.
\item Explicitly state: ``Under the population loss with well-conditioned
      feature covariance $\Sigma \succeq \mu I$\ldots''
\item For LLM claims: drop strong convexity entirely; present as
      ``theory-inspired heuristic.''
\end{itemize}
\end{afterbox}

%% ======================================================================
\section{FATAL Issue \#6: Quantifier Error in Theorem 1}
%% ======================================================================

\begin{beforebox}
\textbf{Theorem 1 (Necessity):}
``\emph{Any} single-timescale algorithm achieving $O(\sqrt{T})$ on stationary
streams suffers $\Omega(T)$ on valid two-scale drift sequences.''
\end{beforebox}

\begin{whywrongbox}
``Single-timescale'' is undefined. The proof sketch only analyzes decaying-LR
SGD with $\eta_t = c/\sqrt{t}$. This excludes:
\begin{itemize}[nosep]
\item Adam/AdaGrad (adaptive preconditioner may enable fast recovery)
\item Restarted OGD / change-point detection algorithms
\item Expert mixtures / strongly adaptive algorithms
\item Constant-LR SGD with periodic averaging (which IS single-model)
\end{itemize}
The quantifier ``any'' is indefensible without a universal information-theoretic
argument.

\textbf{Candidate counterexample}: Adam with constant LR.
After a coefficient jump, the adaptive second-moment estimate is small in the
jump direction, giving an effectively large step size for recovery.
\end{whywrongbox}

\begin{afterbox}
\textbf{Salvageable version:}
\begin{proposition}[Narrow Necessity]
For any algorithm in the class of \emph{single-recursion learners} with
monotonically decaying step size $\eta_t = c \cdot t^{-\alpha}$,
$\alpha \in (0,1]$, on strongly convex batch loss:
if the algorithm achieves $O(\sqrt{T})$ stationary regret,
then there exists a two-rate drift sequence with $K_T = \Theta(\sqrt{T})$
jumps of size $\Delta \geq \Delta_0 > 0$ such that expected dynamic regret
is $\Omega(T)$.
\end{proposition}
\end{afterbox}

%% ======================================================================
\section{FATAL Issue \#15: Regret Bound Mismatch}
%% ======================================================================

\begin{beforebox}
\textbf{Claim (IDEA\_REPORT.md:11, FINAL\_PROPOSAL.md:37):}
``Dual-timescale achieves $O(\sqrt{T})$ regret.''

\textbf{Theorem 2 bound:}
$\mathrm{Reg}_T \leq C_1\sqrt{T(V_f+1)} + C_2\sqrt{T(V_s+1)} + C_3$
\end{beforebox}

\begin{whywrongbox}
Under the claimed regime $K_T = \Theta(\sqrt{T})$ with constant jump size
$\Delta$:
\[
V_f = K_T \cdot \Delta = \Theta(\sqrt{T})
\]
Substituting into Theorem~2:
\[
\mathrm{Reg}_T \leq C_1\sqrt{T \cdot \Theta(\sqrt{T})} + \cdots
= C_1 \cdot \Theta(T^{3/4}) + \cdots
\]
The bound gives $O(T^{3/4})$, \textbf{not} $O(\sqrt{T})$.
The ``$O(\sqrt{T})$'' claim in executive summaries is \textbf{false}
under the theorem's own parameter regime.
\end{whywrongbox}

\begin{afterbox}
\textbf{Options:}
\begin{enumerate}[nosep]
\item Correct the claim to $O(T^{3/4})$ (honest but less impressive)
\item Prove a sharper theorem for jump processes that exploits window-OLS
      exact recovery (would give $O(\sqrt{T} + K_T \cdot w)$ where $w$ is
      window size, which IS $O(\sqrt{T})$ for constant $w$)
\item State the bound for the constant-$V_f$ regime where $O(\sqrt{T})$
      actually holds
\end{enumerate}
\end{afterbox}

%% ======================================================================
\section{FATAL Issues \#18--19: Corollary 1 Failures}
%% ======================================================================

\begin{beforebox}
\textbf{Corollary 1:}
``Fisher-precision averaging $\theta_{\text{merge}} = (\sum F_i)^{-1}(\sum F_i \theta_i)$
decomposes into distillation on recent data + replay on past data.''
\end{beforebox}

\begin{whywrongbox}
\textbf{Two independent fatal problems:}
\begin{enumerate}
\item \textbf{Shape mismatch}: WM-LoRA has rank 16, LT-LoRA has rank 64.
      Direct parameter averaging $\theta_{\text{merge}} = \ldots$ is
      \emph{not even defined}. The code itself states: ``direct parameter
      merge is not possible'' (streaming\_mirror\_lora.py:699).

\item \textbf{No derivation}: The claimed decomposition
      ``Fisher merge = distillation + replay'' has no formal derivation
      anywhere. The code implements iterative KL divergence minimization +
      NTP replay SGD, which is a multi-step optimization procedure, not a
      closed-form precision-weighted average.
\end{enumerate}

\textbf{Empirical contradiction}: Saved E5 results show Fisher-weighted
merge performs \emph{worse} than uniform averaging and EMA.
\end{whywrongbox}

\begin{afterbox}
\begin{remark}[Heuristic Analogy]
Under a Laplace approximation to the posterior over shared-space parameters,
the optimal merge of two Gaussian posteriors
$\mathcal{N}(\theta_i, F_i^{-1})$ is the precision-weighted mean.
In the dual-LoRA architecture with mismatched ranks, this motivates
(but does not formally justify) the KL distillation + replay procedure
as an approximate behavioral merge.
\end{remark}
\end{afterbox}

%% ======================================================================
\section{Counterexample Summary}
%% ======================================================================

\begin{center}
\begin{tabular}{clll}
\toprule
\textbf{CE} & \textbf{Target} & \textbf{Status} & \textbf{Description} \\
\midrule
1 & A1 (strong convexity)      & Verified   & rank-1 Hessian for $d>1$ \\
2 & V\_s gauge invariance       & Verified   & rotation $U \to UR$ changes distance \\
3 & ``Any single-timescale''   & Candidate  & Adam may recover fast from jumps \\
4 & $\Omega(T)$ lower bound    & Verified   & $\Delta_T = T^{-2} \to$ regret $\to 0$ \\
5 & Thm2 sufficiency           & Verified   & E1 saved: dual worse in full drift \\
6 & Cor3 selectivity           & Verified   & E4 saved: selective worse \\
7 & Cor1 Fisher merge          & Verified   & E5 saved: Fisher worse than uniform \\
\bottomrule
\end{tabular}
\end{center}

%% ======================================================================
\section{Proof-Obligation Diff}
%% ======================================================================

\subsection{Before (Current State)}

\begin{itemize}[nosep]
\item 3 theorems claimed --- \textcolor{red}{0 proved}
\item 3 corollaries claimed --- \textcolor{red}{0 derived}
\item 8 assumptions listed --- \textcolor{red}{0 verified for LLM setting}
\item 7 micro-claims identified --- \textcolor{red}{2 with verified counterexamples}
\end{itemize}

\subsection{After (If Option A Salvage Applied)}

\begin{itemize}[nosep]
\item 2 propositions (narrowed) --- provable in linear-Gaussian setting
\item 1 remark (Fisher analogy) --- standard Laplace approximation
\item 1 heuristic ($B^*$) --- from stylized cost model
\item Corollary~3 --- dropped or labeled conjecture
\item LLM claims --- ``theory-inspired,'' not ``theory-derived''
\end{itemize}

%% ======================================================================
\section{Recommended Path Forward}
%% ======================================================================

\begin{insightbox}
\textbf{Recommended: Option A --- Narrow theorem + strong empirical story}

\begin{enumerate}
\item \textbf{Write 2 narrow propositions} for the linear-Gaussian
      setting with expected loss, decaying-step OGD, known subspace.
      Estimated: 2--3 pages of LaTeX proofs.
\item \textbf{Label as ``motivating theory''} (Section 3, 1--2 pages).
      Explicitly state modeling assumptions and note that LLM loss
      violates them.
\item \textbf{Focus the paper on empirical contribution}:
      Mirror-LoRA with Fisher trust-region works well on real LLMs.
      13 methods compared, comprehensive ablations.
\item \textbf{Wire or remove dead code}: Fisher-precision merge,
      invariant detector, and optimal-$B^*$ are never called.
      Either integrate them into the consolidation path or remove
      the corresponding claims.
\item \textbf{Fix synthetic benchmark}: Use $K_T = \Theta(\sqrt{T})$
      for Thm1 validation; fix gauge-dependent metrics; use expected loss.
\end{enumerate}
\end{insightbox}

\begin{warningbox}
\textbf{Do NOT submit with current theorem claims.}
Multiple verified counterexamples from the project's own experiments
directly contradict Corollaries 1--3. Submitting these as formal
results would likely result in desk rejection or strong negative reviews.
\end{warningbox}

%% ======================================================================
\section{Summary: Now Proven / Still Assumed / Open}
%% ======================================================================

\begin{center}
\begin{tabular}{lp{10cm}}
\toprule
\textbf{Category} & \textbf{Items} \\
\midrule
\textcolor{green!60!black}{\textbf{Proven}} &
  (nothing --- no formal proofs exist) \\
\midrule
\textcolor{orange!80!black}{\textbf{Assumed}} &
  Strong convexity (false for LLMs); two-rate decomposition (unverified
  for LLMs); bounded drift; low-rank structure \\
\midrule
\textcolor{red!70!black}{\textbf{Open / Blocked}} &
  Formal proofs for Thm1-2; Fisher merge = distillation derivation;
  invariant coordinate definition; optimal $B^*$ derivation from regret
  bound; extension beyond linear-Gaussian setting \\
\bottomrule
\end{tabular}
\end{center}

\end{document}
